% Options for packages loaded elsewhere
\PassOptionsToPackage{unicode}{hyperref}
\PassOptionsToPackage{hyphens}{url}
\PassOptionsToPackage{dvipsnames,svgnames,x11names}{xcolor}
\documentclass[
  10pt,
  fontset=fandol]{ctexart}
\usepackage{xcolor}
\usepackage[margin=16mm]{geometry}
\usepackage{amsmath,amssymb}
\setcounter{secnumdepth}{-\maxdimen} % remove section numbering
\usepackage{iftex}
\ifPDFTeX
  \usepackage[T1]{fontenc}
  \usepackage[utf8]{inputenc}
  \usepackage{textcomp} % provide euro and other symbols
\else % if luatex or xetex
  \usepackage{unicode-math} % this also loads fontspec
  \defaultfontfeatures{Scale=MatchLowercase}
  \defaultfontfeatures[\rmfamily]{Ligatures=TeX,Scale=1}
\fi
\usepackage{lmodern}
\ifPDFTeX\else
  % xetex/luatex font selection
\fi
% Use upquote if available, for straight quotes in verbatim environments
\IfFileExists{upquote.sty}{\usepackage{upquote}}{}
\IfFileExists{microtype.sty}{% use microtype if available
  \usepackage[]{microtype}
  \UseMicrotypeSet[protrusion]{basicmath} % disable protrusion for tt fonts
}{}
\makeatletter
\@ifundefined{KOMAClassName}{% if non-KOMA class
  \IfFileExists{parskip.sty}{%
    \usepackage{parskip}
  }{% else
    \setlength{\parindent}{0pt}
    \setlength{\parskip}{6pt plus 2pt minus 1pt}}
}{% if KOMA class
  \KOMAoptions{parskip=half}}
\makeatother
\setlength{\emergencystretch}{3em} % prevent overfull lines
\providecommand{\tightlist}{%
  \setlength{\itemsep}{0pt}\setlength{\parskip}{0pt}}
\usepackage{amsmath,amssymb,mathtools,needspace}
\xeCJKDeclareCharClass{CJK}{"2032,"0400->"04FF,"2160->"216B,"2460->"2473,"25A0,"25C7,"25CB,"25CF}
\IfFontExistsTF{Arial Unicode MS}{\xeCJKsetup{AutoFallBack=true}\setCJKfallbackfamilyfont{\CJKrmdefault}{Arial Unicode MS}}{}
\setcounter{secnumdepth}{0}
\setlength{\emergencystretch}{2em}
\allowdisplaybreaks
\usepackage{bookmark}
\IfFileExists{xurl.sty}{\usepackage{xurl}}{} % add URL line breaks if available
\urlstyle{same}
\hypersetup{
  pdftitle={祖冲之``缀术''之谜},
  colorlinks=true,
  linkcolor={Maroon},
  filecolor={Maroon},
  citecolor={Blue},
  urlcolor={Blue},
  pdfcreator={LaTeX via pandoc}}

\title{祖冲之``缀术''之谜}
\author{}
\date{}

\begin{document}
\maketitle

\subsubsection{12.1.2　祖冲之``缀术''之谜}\label{ux7956ux51b2ux4e4bux7f00ux672fux4e4bux8c1c}

祖冲之是公元 5 世纪南北朝时期的数学家。

\begin{quote}
校注（agent 补充）：原页印作``作为正多边形的周长获得圆周率的近似值''，语句疑有排印问题，转录保留原文。

校注（agent 补充）：本页阿尔·卡西圆周率末段印作 \(\ldots792\,23\)，PDF302（书页289）同一结果印作 \(\ldots793\,23\)。独立高精度计算所得 \(\pi=3.141\,592\,653\,589\,793\,238\ldots\)，支持本页``792''为排印错误；此处仍保留原数值。
\end{quote}

\href{../../source-images/pdf-295.jpeg}{原页图像}

祖冲之在数学方面的重大成就，当首推关于圆周率的计算。据《隋书·律历志》记载：

``祖冲之更造密法，以圆径一亿为一丈，圆周盈数三丈一尺四寸一分五厘九毫二秒七忽，肭数三丈一尺四寸一分五厘九毫二秒六忽，正数在盈肭二限之间。''

这就是说，祖冲之定出了圆周率的取值范围为

\[
3.141\,592\,6<\pi<3.141\,592\,7
\]

这个惊人的纪录领先世界一千多年。

此后一直到 15 世纪，再没有出现比这更好的结果。这项辉煌的数学成就，在千年漫长岁月中一直处于世界领先的地位。

祖冲之的 \(\pi\) 值究竟是怎样求出来的呢？据考证祖冲之的算法载于他的《缀术》一书中，可惜该书久已失传，致使这一奇妙算法成了数学史上一桩千古疑案。

后人提出了诸多猜测和假说。

华罗庚先生在《高等数学引论》一书中，以``祖冲之计算圆周率的方法''为题提出了自己的见解，他说：``祖冲之证明了圆周率 \(\pi\) 落在

\[
3.141\,592\,6<\pi<3.141\,592\,7
\]

之间，这是数学史上的一个光辉成就。他的算法也是极限的最好说明！他从单位圆的内接的和外切的正六边形出发，显然圆夹在这两个六边形之间，再作内接的和外切的正 12 边形、正 24 边形、\ldots\ldots、正 \(6\cdot2^{n-1}\) 边形等等，边数愈多，内接的和外切的正 \(6\cdot2^{n-1}\) 边形的面积就愈接近圆的面积，由此可以逐步地精确地算出圆周的长度。''

钱宝琮先生主编的《中国数学史》指出，``缀术失传，祖冲之推算圆周率的方法难以详考。''不过，如果用多边形逼近实现如此高的精度，需要割到 24 576 边形。在筹算的古代，完成如此繁浩的计算量是极为困难的。

一个无可置疑的事实是，祖冲之计算圆周率肯定使用了某种``绝技''。

祖冲之称他的计算技术为``缀术''。据史书记载，缀术``时人称之精妙''，赞扬它``指要精密，算氏之最''。《隋书·律历志》说，祖冲之所著之书名《缀术》，``学官莫能究其深奥''。传说唐代指定《缀术》一书为朝廷钦定的数学教材。

``缀术''是什么？

祖冲之的原著《缀术》已千年失传，无从查察，人们还能还原其``真面目''吗？

\begin{center}\rule{0.5\linewidth}{0.5pt}\end{center}

\subsection{相关原页校注（agent 补充）}\label{ux76f8ux5173ux539fux9875ux6821ux6ce8agent-ux8865ux5145}

以下校注来自本节覆盖的原页，可能也涉及同页相邻小节；原文照录与校正说明分开保留。

\subsubsection{原书 PDF 页 295 的校注}\label{ux539fux4e66-pdf-ux9875-295-ux7684ux6821ux6ce8}

\href{../pages/pdf-295.md}{完整原页转录} · \href{../../source-images/pdf-295.jpeg}{原页图像}

\begin{quote}
校注（agent 补充）：《隋书·律历志》引文中``肭数''\,``盈肭''的``肭''按原页字形保留。
\end{quote}

\subsection{本节覆盖原页的校注记录}\label{ux672cux8282ux8986ux76d6ux539fux9875ux7684ux6821ux6ce8ux8bb0ux5f55}

下列记录按原页关联，含同页相邻小节的校注；计算时同时读取上文校注和完整原页。

\begin{itemize}
\tightlist
\item
  \texttt{CH12-E001}：原书 PDF 页 294
\item
  \texttt{CH12-E002}：原书 PDF 页 294, 302
\end{itemize}

\end{document}
